\section{Energy Eigenvalues Classification}

The energy spectrum of a quantum system typically contains many energy values. The problem then arises of devising a classification method for the energy eigenvalues. There does not exist any a priori ideal classification scheme. However, symmetry provides a natural one. In this section, we show that when the system is characterized by a dynamical symmetry group \(G\), the energy eigenvalues can be efficiently classified based on the unitary representations of \(G\) according to which the energy eigenstates transform under the action of \(G\).

We assume that the system’s Hamiltonian \(H\) has a \textbf{completely discrete spectrum} and that each energy eigenvalue has \textbf{finite degeneracy}. Let \(E\) be an energy eigenvalue of \(H\) and let \(\mathcal{H}_E\) be the corresponding eigenspace, which is \(d_E\)-dimensional given a degeneracy \(d_E\); this subspace is spanned by all vectors \(\Psi\) which respect the eigenvalue equation
\[
    H \Psi = E \Psi.
\]
Let now \(\Psi\) be a vector of \(\mathcal{H}_E\), then for any symmetry transformation \(g\) of \(G\) we have
\[
    H U(g)\Psi = U(g) H \Psi = E U(g) \Psi,
\]
as a direct consequence of \eqref{eq:dynamical_symmetry_hamiltonian_invariance_unitary_representation_definition}. This means that \(U(g)\Psi\) is also an eigenvector of \(H\) with the same eigenvalue \(E\). Thus, the eigenspace \(\mathcal{H}_E\) is invariant under the action of \(G\). The unitary operators \(U(g)\) representing the symmetry transformations of \(G\) on the Hilbert space \(\mathcal{H}\) of the system restrict to unitary operators on the eigenspace \(\mathcal{H}_E\), thus providing a unitary representation of \(G\) on \(\mathcal{H}_E\). The energy eigenvalue \(E\) can be classified according to the unitary representation of \(G\) according to which the energy eigenstates in \(\mathcal{H}_E\) transform under the action of \(G\).

\subsection{Unitary Representations on the Eigenspaces}

If we now consider \( \{ \Phi_i \} \) an ON basis for the subspace \(\mathcal{H}_E\), then these vectors respect the eigenvalue equation and the orthonormality condition
\[
    \begin{dcases}
        H \Phi_i = E \Phi_i, \\
        \bra{\Phi_i} \ket{\Phi_j} = \delta_{ij}.
    \end{dcases}
\]
Since the unitary representation of \(G\) on \(\mathcal{H}_E\) is given by the restriction of the unitary representation of \(G\) on \(\mathcal{H}\), we have that for any symmetry transformation \(g\) of \(G\) and any vector \(\Phi_i\) of the ON basis \(\{\Phi_i\}\) of \(\mathcal{H}_E\), the transformed vector \(U(g)\Phi_i\) is a linear combination of the vectors of the ON basis \(\{\Phi_i\}\):
\begin{equation}
    U(g)\Phi_i = \sum_{j=1}^{d_E} \Phi_j D_{ji}(g),
    \label{eq:unitary_representation_eigenspace_linear_combination_definition}
\end{equation}
where \(D_{ji}(g)\) are the complex, \(g\)-dependent matrix elements of the unitary representation of \(G\) on \(\mathcal{H}_E\). We aim to prove that the family of matrices \(D(g)\) with elements \(D_{ji}(g)\) forms a unitary representation of the symmetry group \(G\) on the eigenspace \(\mathcal{H}_E\).
\begin{proposition}
    For each element \(g\) of the symmetry group \(G\), there exists a \(d_E \times d_E\) complex matrix \(D(g) = (D_{ji}(g))\) which is unitary:
    \begin{equation}
        \begin{aligned}
            D(g \cdot g^{\prime}) & = D(g) D(g^{\prime}),               \\
            D(g)^{\dagger} D(g)   & = D(g) D(g)^{\dagger} = \mathbb{I},
        \end{aligned}
        \label{eq:unitary_representation_eigenspace_definition}
    \end{equation}
    for all \(g\), \(g'\) in \(G\). The family of unitary matrices \(D(g)\) forms a unitary representation of the symmetry group \(G\) on the eigenspace \(\mathcal{H}_E\) corresponding to the energy eigenvalue \(E\).
\end{proposition}

\begin{proof}
    As \(U(g)\) is a unitary representation of \(G\) on the Hilbert space \(\mathcal{H}\) of the system, we know that for any \(g\), \(g'\) in \(G\) we have
    \[
        U(g \cdot g^{\prime}) \Phi_i = U(g) U(g^{\prime}) \Phi_i.
    \]
    Since the vectors \(\Phi_i\) form a basis for the eigenspace \(\mathcal{H}_E\), the above equation implies that the matrix elements of \(U(g)\) in this basis satisfy the same multiplication rule as the matrices \(D(g)\), from \eqref{eq:unitary_representation_eigenspace_linear_combination_definition}:
    \[
        U(g \cdot g^{\prime}) \Phi_i = \sum_{j=1}^{d_E} \Phi_j D_{ji}(g \cdot g^{\prime}),
    \]
    and
    \[
        \begin{aligned}
            U(g) U(g^{\prime}) \Phi_i & = U(g) \sum_{j} \Phi_j D_{ji}(g^{\prime}) = \sum_{j} U(g) \Phi_j D_{ji}(g^{\prime})                                                            \\
                                      & = \sum_{j} \left( \sum_{k} \Phi_k D_{kj}(g) \right) D_{ji}(g^{\prime}) = \sum_{k} \Phi_k \left( \sum_{j} D_{kj}(g) D_{ji}(g^{\prime}) \right).
        \end{aligned}
    \]
    Now if we put together the above two equations, we can prove the first relation in \eqref{eq:unitary_representation_eigenspace_definition}:
    \[
        D_{ji} (g \cdot g^{\prime}) = \sum_{k} D_{kj}(g) D_{ji}(g^{\prime}),
    \]
    which is the multiplication rule for the matrices \(D(g)\).

    The second relation in \eqref{eq:unitary_representation_eigenspace_definition} is a consequence of \eqref{eq:unitary_representation_eigenspace_linear_combination_definition} and the orthonormality of the basis \(\{\Phi_i\}\):
    \[
        \begin{aligned}
            \delta_{ij} & = \bra{\Phi_i} \ket{\Phi_j} = \bra{U(g) \Phi_i} \ket{U(g) \Phi_j}                                                               \\
                        & = \bra{\sum_{k} \Phi_k D_{ki}(g)} \ket{\sum_{l} \Phi_l D_{lj}(g)} = \sum_{k, l} D_{ki}(g)^* \bra{\Phi_k} \ket{\Phi_l} D_{lj}(g) \\
                        & = \sum_{k} D_{ik}(g)^{\dagger} D_{kj}(g) = (D(g)^{\dagger} D(g))_{ij},
        \end{aligned}
    \]
    where in the last steps we have transposed the indices and used the definition of the adjoint matrix \(D(g)^{\dagger}\). This proves that \(D(g)^{\dagger} D(g) = \mathbb{I}\).
\end{proof}

Having proven \eqref{eq:unitary_representation_eigenspace_definition}, we know \(g \mapsto D(g)\) to be a unitary representation of the symmetry group \(G\) on the eigenspace \(\mathcal{H}_E\) corresponding to the energy eigenvalue \(E\), and from \eqref{eq:unitary_representation_eigenspace_linear_combination_definition} we know the basis vectors \(\Phi_i\) to transform under the action of \(G\) through \(D(g)\).

The energy eigenvalue \(E\) can thus be classified according to the unitary representation \(D(g)\), where the energy eigenstates in \(\mathcal{H}_E\) transform under the action of \(G\). The problem is that there are infinitely many orthonormal bases of the energy eigenspace \(\mathcal{H}_E\) of \(E\). The representation \(D(g)\) depends on the chosen basis \(\{\Phi_i\}\). So \(D(g)\) by itself cannot be used to distinguish \(E\) from other energy eigenvalues. Let us analyze this point in greater detail.

\begin{proposition}
    Let \(\{\Phi_{\ i}^1\}\) and \(\{\Phi_{\ i}^2\}\) be two ON bases for the eigenspace \(\mathcal{H}_E\) corresponding to the energy eigenvalue \(E\). Then, let the corresponding unitary representations be \(D^1(g)\) and \(D^2(g)\), defined by \eqref{eq:unitary_representation_eigenspace_linear_combination_definition} with respect to the bases \(\{\Phi_{\ i}^1\}\) and \(\{\Phi_{\ i}^2\}\), respectively.

    There exist a \(d_E \times d_E\) unitary matrix \(A\) such that its elements \(A_{ij}\) relate the two bases \(\{\Phi_{\ i}^1\}\) and \(\{\Phi_{\ i}^2\}\) as
    \begin{equation}
        \Phi_{\ i}^2 = \sum_{j=1}^{d_E} \Phi_{\ j}^1 A_{ji},
        \label{eq:unitary_representation_eigenspace_basis_change_definition}
    \end{equation}
    and the two unitary representations \(D^1(g)\) and \(D^2(g)\) are related by the similarity transformation
    \begin{equation}
        D^2(g) = A^{-1} D^1(g) A,
        \label{eq:unitary_representation_eigenspace_similarity_transformation_definition}
    \end{equation}
    for all \(g\) in \(G\). Thus, the two unitary representations \(D^1(g)\) and \(D^2(g)\) are equivalent representations of the symmetry group \(G\).
\end{proposition}

\begin{proof}
    Using \eqref{eq:unitary_representation_eigenspace_linear_combination_definition}, we can develop \eqref{eq:unitary_representation_eigenspace_basis_change_definition} as
    \[
        \begin{aligned}
            U(g) \Phi_{\ i}^2 & = U(g) \sum_{j=1}^{d_E} \Phi_{\ j}^1 A_{ji} = \sum_{j=1}^{d_E} U(g) \Phi_{\ j}^1 A_{ji} = \sum_{j=1}^{d_E} \left( \sum_{k=1}^{d_E} \Phi_{\ k}^1 D_{kj}^1(g) \right) A_{ji} \\
                              & = \sum_{k=1}^{d_E} \Phi_{\ k}^1 \left( \sum_{j=1}^{d_E} D_{kj}^1(g) A_{ji} \right) = \sum_{k=1}^{d_E} \Phi_{\ k}^1 (D^1(g) A)_{ki},
        \end{aligned}
    \]
    on one hand, and
    \[
        \begin{aligned}
            U(g) \Phi_{\ i}^2 & = \sum_{k=1}^{d_E} \Phi_{\ k}^2 D_{ki}^2(g) = \sum_{k=1}^{d_E} \left( \sum_{j=1}^{d_E} \Phi_{\ j}^1 A_{jk} \right) D_{ki}^2(g)      \\
                              & = \sum_{j=1}^{d_E} \Phi_{\ j}^1 \left( \sum_{k=1}^{d_E} A_{jk} D_{ki}^2(g) \right) = \sum_{j=1}^{d_E} \Phi_{\ j}^1 (A D^2(g))_{ji}.
        \end{aligned}
    \]
    Putting together the above two equations, we have
    \[
        (D^1(g) A)_{ki} = (A D^2(g))_{ki}, \quad \text{i.e., } D^1(g) A = A D^2(g),
    \]
    which shows that the two unitary representations \(D^1(g)\) and \(D^2(g)\) are related by the similarity transformation \(D^2(g) = A^{-1} D^1(g) A\). This proves that the two unitary representations \(D^1(g)\) and \(D^2(g)\) are equivalent representations of the symmetry group \(G\).
\end{proof}

Thus two unitary representations \(D^1(g)\) and \(D^2(g)\) of the symmetry group \(G\) of the same dimension \(d_E\), related as in \eqref{eq:unitary_representation_eigenspace_similarity_transformation_definition}, for some unitary matrix \(A\), are said to be \textbf{unitarily equivalent}. Thus the unitary representation \(D(g)\) in \eqref{eq:unitary_representation_eigenspace_linear_combination_definition} is independent of the choice of the ON basis \(\{\Phi_i\}\) for the eigenspace \(\mathcal{H}_E\) up to unitary equivalence.

We can in this way characterize the energy eigenvalue \(E\) by a \textit{unitary equivalence class of unitary representations}\footnote{It is a class of equivalence of unitary representations, where the representations are unitarily equivalent.} of \(G\) and use this class to group \(E\) with other eigenvalues sharing the very same class. In this way, \textbf{a symmetry based classification of the energy spectrum is achieved}. Since however there are infinitely many representations of any given group, the classification scheme obtained may not be so useful. Fortunately, generic representations can be decomposed as combinations of a smaller set of more basic representations, which can then be used to refine the scheme.

\subsection{Reducibility and Irreducibility}

Suppose that we are able to chose the orthonormal basis \(\{\Phi_i\}\) of the energy eigenspace \(\mathcal{H}_E\) to be of indexed \(\Phi_{\alpha r}\) by two indices \(\alpha\), \(r\) in such a way that \eqref{eq:unitary_representation_eigenspace_linear_combination_definition} takes the form
\[
    U(g) \Phi_{\alpha r} = \sum_{s} \Phi_{\alpha s} D_{\alpha s r}(g),
\]
for every element \(g\) of \(G\), for certain \(g\)-dependent coefficients \(D_{\alpha s r}(g)\). Therefore, the vectors \(\Phi_{\alpha r}\) for fixed \(\alpha\) mix among themselves under the action of \(G\). This means that, for fixed \(g\), the matrix \(D(g)\) has the block diagonal form
\[
    D(g) = \begin{pmatrix}
        D_1(g) & 0      & \cdots & 0      \\
        0      & D_2(g) & \cdots & 0      \\
        \vdots & \vdots & \ddots & \vdots \\
        0      & 0      & \cdots & D_a(g)
    \end{pmatrix},
\]
and let us assume that the blocks \(D_{\alpha}(g)\) are of dimension \(d_{\alpha}\), which is the smallest non trivial dimension for each block (i.e. the blocks cannot be further decomposed into diagonal blocks by a furhter change of basis \(\Phi_{\alpha r}\)). Then, the unitary representation \(g \mapsto D(g)\) of \(G\) on \(\mathcal{H}_E\) is said to be \textbf{reducible}, and its blocks \(D_{\alpha}(g)\) are said to define an \textbf{irreducible} representations \(g \mapsto D_{\alpha}(g)\) of \(G\) (\textit{irreducible components}).

The energy eigenvalue \(E\) can then be classified through the irreducible representations \(D_{\alpha}(g)\) of \(G\). The former way of proceeding provides a classification scheme of the energy spectrum relying on a more restricted set of elements and so potentially more useful than that furnished by the latter.

We can resume the results obtained in this section as follows:
\begin{itemize}
    \item Each energy eigenvalue \(E\) is characterized by a set of irreducible unitary representations \(D_{\alpha}\) of \(G\), unique up to unitary equivalence, describing how the eigenvectors belonging to \(E\) transform under the action of \(G\).
    \item The energy eigenvalues can be grouped in subsets of eigenvalues sharing the same representation content \(D_{\alpha}\) providing a classification scheme of the energy spectrum.
\end{itemize}

Let us illustrate the above results with a few examples.

\begin{example}[The One-Dimensional Harmonic Oscillator]
    The parity transformation \(P\) is space reflection about the origin. Observation show that the parity group \(C_2 = \{\mathbb{I},P\}\) is a dynamical symmetry group of the one–dimensional harmonic oscillator.

    As well known, the energy eigenvalues of the one–dimensional harmonic oscillator are given by
    \[
        H \Phi_n = E_n \Phi_n. \quad E_n = \hbar \omega (n + 1/2),
    \]
    with \(n = 0, 1, 2, \ldots\), and are non degenerate. The energy eigenstates of the one–dimensional harmonic oscillator transform under the action of the parity transformation \(P\) as
    \[
        U(P) \Phi_n = \Phi_n (-1)^n.
    \]
    Thus the parity group \(C_2\) has two irreducible unitary representations: the trivial representation \(D_+(P) = 1\) and the sign representation \(D_-(P) = -1\). The energy eigenvalues of the one–dimensional harmonic oscillator are classified according to these two irreducible unitary representations of \(C_2\): the energy eigenvalues with even \(n\) transform according to the trivial representation \(D_+\), while those with odd \(n\) transform according to the sign representation \(D_-\). Thus we classify them as even or odd energy eigenvalues according to the parity of \(n\).
\end{example}

\begin{example}[The Hydrogen Atom]
    It can be shown that the proper rotation group \(\mathrm{SO}(3)\) is a dynamical symmetry group of the hydrogen atom. The energy eigenvalues of the hydrogen atom are known to be given by
    \[
        H \Phi_{n l m} = E_n \Phi_{n l m}, \quad E_n = -\frac{m c^2 \alpha^2}{2 n^2}.
    \]
    The energy eigenvalues of the hydrogen atom are degenerate, with degeneracy \(d_{E_n} = n^2\):
    \begin{enumerate}
        \item \(n = 0, 1, 2, \ldots\),
        \item \(l = 0, 1, \ldots, n-1\),
        \item \(m = -l, -l+1, \ldots, l-1, l\),
    \end{enumerate}
\end{example}
and the eigenstates \(\Phi_{n l m}\) are also eigenstates of the total angular momentum operator \(L^2\) and of the component along the quantization axis of the angular momentum operator \(L_z\):
\[
    \begin{dcases}
        L^2 \Phi_{n l m} = \hbar^2 l(l+1) \Phi_{n l m}, \\
        L_0 \Phi_{n l m} = \hbar m \Phi_{n l m}.
    \end{dcases}
\]
Furthermore, the other spherical coordinates of the angular momentum \(L_{\pm}\) act on the eigenstates \(\Phi_{n l m}\) as
\[
    L_{\pm} \Phi_{n l m} = \hbar \sqrt{l(l+1) - m(m \pm 1)} \Phi_{n l (m \pm 1)},
\]
thus acting as ladder operators on the index \(m\).

Since \(\mathbf{L}\) is the \textbf{infinitesimal generator} of the proper rotation group \(\mathrm{SO}(3)\) (which will be discussed in the next section \ref{sec:lie}), the energy eigenstates \(\Phi_{n l m}\) transform under the action of \(\mathrm{SO}(3)\) according to the irreducible unitary representation \(D_l(\mathbf{R})\) of \(\mathrm{SO}(3)\), of dimension \(d_l = 2l + 1\):
\[
    U(\mathbf{R}) \Phi_{n l m} = \sum_{m'=-l}^{l} \Phi_{n l m'} D_{l m' m}(\mathbf{R}),
\]
so that \(D_l(\mathbf{R})\) is \((2l+1)\times(2l+1)\) dimensional, and the mappings \(\mathbf{R} \mapsto D_l(\mathbf{R})\) can be shown to be unitary irreducible representations of \(\mathrm{SO}(3)\). Thus for each energy eigenvalue \(E_n\) of the hydrogen atom, the energy eigenstates \(\Phi_{n l m}\) transform under the action of \(\mathrm{SO}(3)\) according to the \(n\) irreducible \((2l+1)\)-dimensional unitary representation \(D_{l=0,1,\ldots,n-1}\) of \(\mathrm{SO}(3)\). In this way, the symmetry based classification of the energy spectrum is achieved for the hydrogen atom.

\begin{example}[The Sulfur Trioxide Molecule SO\(_3\)]\TODO{put image of the molecule}
    The molecule of sulfur trioxide SO\(_3\) has the shape of a planar equilateral triangle, with the sulfur atom at the center and the three oxygen atoms at the vertices. The symmetry group of this molecule is the dihedral group \(D_3\), which is a subgroup of \(\mathrm{SO}(3)\).

    A planar equilateral triangle has the following symmetries:
    \begin{itemize}
        \item Three rotations of \(120^{\circ}\) about the axis perpendicular to the plane of the triangle, which we call \(C_3\) rotations:
              \[
                  0^{\circ} : \begin{pmatrix}
                      123        \\
                      \downarrow \\
                      123
                  \end{pmatrix}, \quad 120^{\circ} : \begin{pmatrix}
                      123        \\
                      \downarrow \\
                      231
                  \end{pmatrix}, \quad 240^{\circ} : \begin{pmatrix}
                      123        \\
                      \downarrow \\
                      312
                  \end{pmatrix}.
              \]
        \item Three reflections about the axes of symmetry of the triangle, which we call \(\sigma_v\) reflections:
              \[
                  \sigma_v^1 : \begin{pmatrix}
                      123        \\
                      \downarrow \\
                      132
                  \end{pmatrix}, \quad \sigma_v^2 : \begin{pmatrix}
                      123        \\
                      \downarrow \\
                      321
                  \end{pmatrix}, \quad \sigma_v^3 : \begin{pmatrix}
                      123        \\
                      \downarrow \\
                      213
                  \end{pmatrix}.
              \]
    \end{itemize}
    Putting together the above symmetries, we find a total of six symmetries, which form the dihedral group \(D_3\), which coincides with the symmetric group \(S_3\) of permutations of three objects. The energy eigenvalues of the sulfur trioxide molecule SO\(_3\) can be classified according to the irreducible unitary representations of \(D_3\).

    This group has three \textit{unitary equivalence classes of irreducible unitary representations}: the trivial representation \(D_+\), the sign representation \(D_-\) and a two-dimensional representation \(D_2\). They are of dimensions \(d_+ = 1\), \(d_- = 1\) and \(d_2 = 2\), respectively, and we can deduce their dimensionality through the following theorem:
    \begin{theorem}
        The sum of the squares of the dimensions \(d_{\alpha}\) of the irreducible unitary representations \(D_{\alpha}\) of a finite non-abelian group \(G\) is equal to the order \(|G|\) (the number of elements) of the group:
        \[
            \sum_{\alpha} d_{\alpha}^2 = |G|.
        \]
    \end{theorem}
    In our case the only way to satisfy the above relation is to have \(d_+^2 + d_-^2 + d_2^2 = 1^2 + 1^2 + 2^2 = 6\), which is the order of \(D_3\). The energy eigenvalues of the sulfur trioxide molecule SO\(_3\) can be classified according to these three irreducible unitary representations of \(D_3\).
\end{example}